\subsection*{Предметная область}

Задача о наибольшей общей подпоследовательности (Longest Common Subsequence, LCS) --- одна из фундаментальных задач алгоритмики и комбинаторики на словах. Классический алгоритм Вагнера--Фишера~\cite{wagner1974} вычисляет длину LCS двух строк длины $n$ за время $O(n^2)$. Помимо непосредственных приложений (системы контроля версий, утилита \texttt{diff}, выравнивание биологических последовательностей), LCS играет роль модельной задачи в теории случайных процессов на словах.

Пусть $A = a_1 a_2 \ldots a_n$ и $B = b_1 b_2 \ldots b_n$ --- независимые случайные строки, каждая буква которых выбирается равномерно и независимо из алфавита размера $\sigma$. Пусть $L_n$ --- длина LCS этих строк. В 1975~году Чватал и Санкофф~\cite{chvatal1975} доказали, что существует предел
\[
    \gamma_\sigma := \lim_{n \to \infty} \frac{\mathbb{E}[L_n]}{n}.
\]
Эта величина называется \emph{константой Чватала--Санкоффа}. Точное значение $\gamma_\sigma$ не известно ни при каком $\sigma$ и является классической открытой проблемой~\cite{lueker2009, heineman2024}.

\subsection*{Объект и предмет исследования}

\textbf{Объект исследования} --- случайная величина $L_n$ длины LCS двух независимых равномерно случайных строк длины $n$ над алфавитом размера~$\sigma$.

\textbf{Предмет исследования} --- предел $\gamma_\sigma = \lim_{n \to \infty} \mathbb{E}[L_n]/n$ и строгие нижние оценки этой величины, получаемые конечно\-комбинаторными методами.

\subsection*{Актуальность}

Для бинарного алфавита $\sigma = 2$ известные строгие оценки таковы:
\[
    0{,}7927 \leq \gamma_2 \leq 0{,}8263,
\]
где нижняя оценка получена Heineman et al. в~\cite{heineman2024} (2024), верхняя --- Lueker'ом в~\cite{lueker2009} (2009). Эмпирические оценки методом Monte-Carlo показывают $\gamma_2 \approx 0{,}811$, что указывает на существенный разрыв нижней оценки с истинным значением. С 2003 по 2024~год нижняя оценка $\gamma_2$ улучшалась медленно: суммарный прирост за два десятилетия составил всего $\approx 0{,}005$. Каждое новое улучшение, превышающее $10^{-3}$, представляет самостоятельный интерес.

Аналогичный разрыв имеет место и для $\sigma \geq 3$. Существующие подходы (конечные автоматы, расширенные пространства состояний) упираются в экспоненциальный рост пространства состояний, что мотивирует поиск качественно новых вычислительных схем.

Принципиальное преимущество предлагаемого в работе подхода --- разделение двух параметров: размера окна обзора $N$ и максимального сдвига $\mathrm{MAX\_SH}$. Это позволяет наращивать выразительность стратегии (рост $N$) и допустимый «зазор» между указателями (рост $\mathrm{MAX\_SH}$) независимо, удерживая память в управляемых пределах $O(\mathrm{MAX\_SH} \cdot \sigma^{2N})$, тогда как в схемах с единым параметром пространство состояний растёт как $\sigma^{O(N^2)}$ или быстрее.

\subsection*{Цель и задачи}

\textbf{Цель работы} --- получить строгую нижнюю оценку $\gamma_\sigma$, превосходящую известные результаты, для $\sigma \in \{2, 3, \ldots, 10\}$.

\textbf{Задачи:}
\begin{enumerate}
    \item Сформулировать оконную стратегию построения общей подпоследовательности двух случайных строк и доказать, что её ожидаемая нормированная длина является нижней оценкой $\gamma_\sigma$.
    \item Разработать метод динамического программирования, позволяющий точно вычислить оптимальную ожидаемую длину в классе оконных стратегий.
    \item Доказать корректность ДП-уравнений и оценить асимптотическую сложность алгоритма по памяти и времени.
    \item Реализовать алгоритм на C++ с современными оптимизациями (одинарная точность, AVX2, развёртка циклов).
    \item Провести масштабные численные эксперименты для $\sigma \in \{2, \ldots, 10\}$.
    \item Независимо верифицировать результаты методом Monte-Carlo.
    \item Сравнить полученные оценки с лучшими опубликованными результатами.
\end{enumerate}

\subsection*{Методы исследования}

В работе используются: формальное доказательство корректности рекуррентных уравнений ДП индукцией по числу шагов; лемма Фекете~\cite{fekete1923} о супераддитивных последовательностях для обоснования существования предела $f_K/K$; численные эксперименты на языке C++ с использованием векторных расширений AVX2; независимая проверка методом Monte-Carlo по классическому ДП-алгоритму LCS на $500$ парах случайных строк длины до $20\,000$ для каждого $\sigma$.

\subsection*{Основные результаты}

\begin{itemize}
    \item Получена рекордная нижняя оценка $\gamma_2 \geq 0{,}7964$ (предыдущий рекорд $0{,}79266$~\cite{heineman2024}).
    \item Получены рекордные нижние оценки $\gamma_\sigma$ для всех $\sigma \in \{3, \ldots, 10\}$ с приростом от $+0{,}17\,\%$ до $+5{,}42\,\%$ относительно лучших известных результатов.
    \item Все оценки независимо подтверждены методом Monte-Carlo.
\end{itemize}

\subsection*{Структура работы}

Работа состоит из пяти разделов. В разделе~1 даётся обзор предметной области и сводка известных результатов. В разделе~2 формулируется метод оконного ДП. В разделе~3 доказывается корректность метода и анализируется его сложность. Раздел~4 описывает реализацию и историю оптимизаций. Раздел~5 содержит результаты численных экспериментов и Monte-Carlo верификацию.
